\documentclass[11pt]{article}
\usepackage[margin=0.85in]{geometry}
\usepackage{amsmath,amssymb,booktabs,array}
\usepackage[hidelinks]{hyperref}
\setlength{\parindent}{0pt}
\setlength{\parskip}{0.4em}
\setlength{\emergencystretch}{2em}
\newcommand{\ve}{\mathrm{veh\mbox{-}eq}}
\newcommand{\dd}{\mathrm{d}}
\title{Elementary VSL Lane-Drop Opening\\[0.3em]
\large Inherited S13 kernel, complete time and physical sensitivity}
\author{Method-first task-definition review}
\date{12 October 2026}
\begin{document}
\maketitle
\vspace{-1em}

\textbf{Status.} Opening support only. These equations restate the accepted
S13 conditional member in SI units; they do not qualify an applied model,
select an evaluator or authorize implementation.

\section*{1. Traffic meaning, geometry and state}
At a permanent three-to-two lane drop, slower common speed costs travel time
but can leave less merging for the narrow exit. Does initial preparation change
the useful speed perturbation, and does one continuous statistic suffice?

A fixed cohort of mass $M>0$ is at the entrance of a preparation zone of length
$L>0$. The ending lane is explicit; the two continuing lanes form a symmetric
group. For fixed preparation $s$, its initial ending-lane fraction is $p_s$:
\begin{equation}
 e_s(0)=p_s M,\qquad c_s(0)=(1-p_s)M,\qquad 0<p_s<1.
 \label{eq:initial}
\end{equation}
Here $e_s,c_s,M$ are continuous $\ve$; $L$ is metres; time is seconds.
Lane labels, topology, routes and preparation are fixed during differentiation.
Predecision history equality is conditional on S05/S13, not demonstrated in an
independent evaluator.

The only action is constant common advice $v\in[v_{\min},v_{\max}]$, in
metres per second. In the admitted ideal longitudinal-response branch,
\begin{equation}
 \tau(v)=\frac{L}{v},\qquad
 \tau_{\min}=\frac{L}{v_{\max}},\quad
 \tau_{\max}=\frac{L}{v_{\min}},\qquad 0<v_{\min}<v_{\max}.
 \label{eq:traversal}
\end{equation}
Advice equals realized longitudinal speed only under this assumption. Safety,
congestion and heterogeneous compliance are not modeled by changing the name
of $v$.

\section*{2. Conservative preparation}
The inherited fixed active branch uses $\lambda>0$ in $\mathrm{s}^{-1}$,
with preparation-zone target receiving assumed nonbinding:
\begin{equation}
 \dot e_s=-\lambda e_s,\qquad \dot c_s=\lambda e_s,
 \qquad e_s(t)+c_s(t)=M.
 \label{eq:transfer}
\end{equation}
Consequently, at the choke,
\begin{equation}
 E_s(v):=e_s(\tau(v))=p_sM\exp(-\lambda L/v),\qquad
 c_s(\tau)=M-E_s(v).
 \label{eq:remaining}
\end{equation}
This exponential is a bounded analytical adaptation of conservative mandatory
transfer, not a calibrated empirical lane-changing law. The longitudinal batch
and lateral bookkeeping conserve continuous traffic mass; no rounding or
vehicle identity is used.

\newpage
\section*{3. Finite receiving and complete time}
Let ordinary continuing mass require one unit of receiving work per unit mass,
and unfinished merging require $\alpha>1$ units. With common receiving-work
rate $\mu>0$, in normalized $\ve/\mathrm{s}$, the work and clearance duration are
\begin{equation}
 W_s(v)=M+(\alpha-1)E_s(v),\qquad T_s(v)=W_s(v)/\mu.
 \label{eq:work}
\end{equation}
The factor $\alpha$ changes receiving work, not vehicle mass. S13 assumes
homogeneous proportional batch service: for elapsed service time $u$,
\begin{equation}
 D_s(u)=M u/T_s,\quad 0\le u\le T_s,\qquad D_s(T_s)=M.
 \label{eq:departures}
\end{equation}
All mass spends $\tau$ in preparation; thereafter the remaining population
declines linearly to zero. Complete total system time is therefore
\begin{align}
 F_s(v)&=M\tau+\int_0^{T_s}(M-D_s(u))\,\dd u
       =M\tau+\frac{M T_s}{2}, \label{eq:total}\\
 J_s(\tau)&:=F_s/M
       =\tau+\frac{M}{2\mu}+Kp_s e^{-\lambda\tau},
       \qquad K=\frac{(\alpha-1)M}{2\mu}. \label{eq:mean}
\end{align}
$F_s$ has units $\ve\,\mathrm{s}$; the mean $J_s$ has units seconds,
not seconds per $\ve$. Fixed $M$ makes their action ranking identical.
Common action-independent prefix/downstream costs do not change that ranking;
any action-dependent waiting or travel must be charged. A different service
priority cannot inherit $T_s/2$ merely because final clearance work is unchanged.

\section*{4. Physical sensitivity and inherited open regime}
Within the fixed smooth branch, the inherited derivatives are
\begin{align}
 \frac{\dd J_s}{\dd\tau}&=1-K\lambda p_s e^{-\lambda\tau},\qquad
 \frac{\dd^2J_s}{\dd\tau^2}=K\lambda^2p_s e^{-\lambda\tau}>0,\label{eq:taugrad}\\
 \frac{\dd F_s}{\dd v}&=-\frac{ML}{v^2}
       \left(1-K\lambda p_s e^{-\lambda L/v}\right).
 \label{eq:speedgrad}
\end{align}
The last derivative has units $\ve\,\mathrm{s}^2/\mathrm{m}$. It is the
physical common-speed sensitivity; a traversal-time derivative is an internal
coordinate and must be mapped back by $\dd\tau/\dd v=-L/v^2$.

For $0<p_L<p_H<1$, the accepted certificate assumes the strict inequalities
\begin{equation}
 K\lambda p_L<e^{\lambda\tau_{\min}}
 <K\lambda p_H<e^{\lambda\tau_{\max}}.
 \label{eq:regime}
\end{equation}
Then the low-ending-lane state uniquely selects $\tau_{\min}$, whereas the high
state selects the interior $\tau_H^*=\lambda^{-1}\log(K\lambda p_H)$.
At $v_{\max}$ the feasible perturbation is a speed reduction: it improves the
high state and worsens the low state. The low state is already at its optimum;
an increase above the bound is not a permissible intervention. Strict openness
does not establish that realistic traffic commonly occupies this regime.

\newpage
\section*{5. Minimum information and fair prediction}
The strong aggregate observer is the complete declared predecision history of
aggregate density, flow, speed and cumulative counts. Its exact equality across
preparations is conditional on S05/S13's upstream construction. Equality of a
snapshot alone is insufficient.

Use three roles: aggregate history alone; that history plus initial ending-lane
mass $E_{0,s}=p_sM$; and explicit ending/continuing groups. The scalar model
predicts independently under the inherited law:
\begin{equation}
 E_s(v)=E_{0,s}e^{-\lambda L/v},\qquad c_s(\tau)=M-E_s(v).
 \label{eq:compression}
\end{equation}
Thus it exactly reproduces the grouped objective and derivative for this member.
It receives no future reference correction. Total mass and one continuous
initial merge statistic suffice here; full lane state, order and identity have
not been shown necessary. Comparing aggregate-only with scalar-augmented inputs
is an information test; comparing scalar and grouped inputs is a compression
test at equal information.

\section*{6. Evidence sequence and permitted conclusions}
Conservation, finite receiving, complete accounting and compact equivalence
come first. A later authorized implementation would compare the analytical
derivative with AD and central differences in double precision, using feasible
interior steps with a stable event signature. Endpoint or event cases require
separate one-sided/piecewise interpretation, without smoothing them away.

An independent evaluator must then test common physical perturbations selected
without its outcomes. Correct internal derivatives and internal objective
improvement do not establish external usefulness. Reproducing these equations
with an independent transcription verifies the equations, not empirical traffic
validity. No evaluator is admitted here.

The initial deletion audit removes action-sensitive transfer, removes extra
merge work, inspects nonbinding receiving and removes the hidden-state contrast.
No new deletion or alternative service law is frozen by this note. The complete
time objective stays fixed; changing it to obtain priority invariance would be
a new methodological choice. Strong one-dimensional direct search and the
analytical optimum remain optimizer references; gradients need not be necessary
to solve this scalar task.

\section*{7. Source lineage and scope}
Per-lane continuous conservation and mandatory-transfer structure follow the
accepted FB0030/S12 source audit. Finite receiving, gap-weighted work and
allocation alternatives are grounded in
\href{https://doi.org/10.1007/s11067-013-9193-7}{Carey, Balijepalli and Watling}.
Common VSL at a permanent drop has task precedent in
\href{https://doi.org/10.1080/23249935.2022.2033878}{Zhang et al.}
These sources do not supply the joint exponential/batch model as a calibrated
law. The S13 adaptations and deductions remain explicitly conditional.

This note preserves the accepted S13 result and does not broaden its claim to
family robustness, native trajectory reachability, implementation derivatives,
empirical prevalence, independent usefulness or novelty. The urban connection
is the information--service--sensitivity--intervention evidence chain; its
discrete-order hypothesis remains distinct and unchanged.

\end{document}
